\begin{textsample}{POS Dim 1 – 2011\_c4 – Score -4.00 – 2011\_c4/enade\_2011\_filo\_74486d8bb9.txt}  \label{ex:f1_pos_330}
CAPÍTULO 4 PERCEPÇÃO DA PROVA As análises feitas neste capítulo tratam das percepções dos concluintes da Área de Filosofia sobre a prova aplicada no ENADE/2011.
Estas percepções foram mensuradas por meio de nove questões que avaliaram desde o grau de dificuldade da prova até o tempo gasto para concluí -la.
As percepções sobre a prova foram relacionadas com o desempenho dos estudantes e com a Grande Região de funcionamento do curso.
O questionário de percepção da prova encontra -se ao final do Anexo IV, que traz a reprodução da prova.
O desempenho dos estudantes foi classificado em quatro quartos.
Para tanto, esse desempenho foi ordenado de forma ascendente.
O percentil 25, P25, também conhecido como primeiro quartil, é a nota de desempenho que deixa um quarto ( 25 \% ) dos valores observados abaixo e três quartos acima.
A Figura 1 apresenta uma ilustração deste conceito.
O quarto inferior de desempenho é composto pelas notas abaixo do primeiro quartil.
Já o percentil 75, P75, também conhecido como terceiro quartil, é o valor para o qual há três quartos ( 75 \% ) dos dados abaixo e um quarto acima dele.
O quarto superior de desempenho é composto pelas notas iguais ou acima do terceiro quartil.
O percentil 50, P50, também conhecido como mediana, é o valor que divide as notas em dois conjuntos de igual tamanho.
O segundo quarto inclui valores entre o primeiro quartil ( P25 ) e a mediana.
O terceiro quarto contém os valores entre a mediana ( P50 ) e o terceiro quartil ( P75 ).
Vale ressaltar que percentis, quartis e medianas são pontos que não obrigatoriamente pertencem ao conjunto original de dados, ao passo que os quartos são subconjuntos dos dados originais.
A seguir, serão apresentados gráficos com resultados selecionados, relativos às nove questões avaliadas por grupos de estudantes.
Os gráficos apresentam nas barras o percentual de alunos que assinalaram uma das opções ou a soma das porcentagens daqueles que assinalaram duas ( ou três ) delas.
Por exemplo, para as questões 1 e 2, os gráficos apresentam a porcentagem total de participantes que assinalaram as opções ( D ) difícil e ( E ) muito difícil.
Em cada barra foram assinalados também os extremos do intervalo de confiança de 95 \% como linhas verticais unidas por uma linha horizontal na forma da letra H maiúscula, semelhantemente aos gráficos do Capítulo 3.
As Tabelas no Anexo II apresentam os valores absolutos e a distribuição percentual das alternativas válidas das nove questões, segundo o mesmo recorte de desempenho dos alunos e Grande Região de funcionamento do curso. 4.1 GRAU DE DIFICULDADE DA PROVA 4.1.1 Componente de Formação Geral Ao avaliarem “ Qual o grau de dificuldade desta prova na parte de Formação Geral? ” ( Questão 1 ), 26,7 \% do grupo de inscritos e presentes optaram pelas alternativas difícil ou muito difícil.
Entretanto, para mais da metade dos estudantes ( 55,3, \% ), o Componente de Formação Geral da prova foi considerado com grau de dificuldade médio ( Gráfico 4.1, Gráfico 4.2 e, no Anexo II, a Tabela II.1 ).
O percentual de estudantes que consideraram a prova como difícil ou muito difícil foi maior na região Centro-Oeste, onde a proporção foi de 29,8 \%, enquanto a de menor incidência foi a Sudeste, com 24,9 \%.
No Gráfico 4.1 é possível observar que esta diferença não é estatisticamente significativa.
Nas Grandes Regiões, a proporção de presentes à prova que consideraram o Componente de Formação Geral como sendo de grau de dificuldade médio esteve entre 53,0 \% na região Sul e 58,8 \% na região Norte.
O percentual de alunos que consideraram a prova difícil ou muito difícil apresentou uma tendência decrescente com o aumento do desempenho, variando de 36,6 \% ( 1º quarto ) até 18,2 \% ( 4º quarto ).
A alternativa modal para esta pergunta foi médio, com 49,1 \% no quarto inferior e com os valores nos outros três quartos de desempenho ficando em torno de 57 \% com menos de 2 \% de diferença entre o maior e o menor percentual. 4.1.2 Componente de Conhecimento Específico Ao responderem à Questão 2 – “ Qual o grau de dificuldade desta prova na parte de Componente Específico? ” – 39,9 \% do grupo de estudantes classificaram -na como difícil ou muito difícil.
Além disso, o Componente de Conhecimento Específico da prova foi considerado com grau de dificuldade médio por 51,4 \% dos alunos ( Gráfico 4.3, Gráfico 4.4, e, no Anexo II, a Tabela II.2 ).
A análise das respostas dos estudantes quanto ao grau de dificuldade do Componente de Conhecimento Específico da prova, agregado por Grande Região, mostra que a diferença entre a maior e a menor proporção de alunos que a avaliaram como difícil ou muito difícil é estatisticamente significativa : 33,2 \% na região Norte e 45,7 \% na Sul.
O percentual de alunos que classificaram o grau de dificuldade como médio, no Componente de Conhecimento Específico, variou de 45,9 \% a 57,9 \%, para as mesmas regiões, Sul e Norte, respectivamente.
Considerando -se a avaliação da dificuldade das questões do Componente de Conhecimento Específico da prova, de acordo com o desempenho dos estudantes, não se observa diferença estatisticamente significativa.
Em todos os quartos, a proporção dos que classificaram a parte específica como difícil ou muito difícil variou de 42,1 \% ( 4º quarto ) a 38,6 \% ( 3º quarto ).
A alternativa modal para a Questão 2 foi o médio, com 49,6 \% do quarto inferior e 49,4 \% do superior optando por esta resposta ; o percentual mais elevado foi no 2º quarto ( 53,4 \% ). 4.2 EXTENSÃO DA PROVA EM RELAÇÃO AO TEMPO TOTAL Indagados quanto à extensão da prova, em relação ao tempo total oferecido para a sua resolução ( Questão 3 ), os estudantes apontaram, com maior incidência, a alternativa que considerava a extensão adequada, para todas as agregações consideradas ( Gráfico 4.5, Gráfico 4.6, e, no Anexo II, a Tabela II.3 ).
O percentual de alunos que responderam ser a extensão da prova adequada foi de 47,2 \%.
Já 46,7 \% dos inscritos presentes consideraram que o exame foi longo ou muito longo e menos do que 7 \% o avaliaram como curto ou muito curto.
Entre as Grandes Regiões a proporção daqueles que avaliaram a prova como longa ou muito longa em relação ao tempo total destinado à sua resolução variou pouco : de 41,9 \% na região Norte até 50,1 \% na região Nordeste.
A diferença entre as regiões Norte e Nordeste não é estatisticamente significativa.
Considerando -se o desempenho dos alunos, optaram pela extensão da prova adequada percentuais próximos que ficaram em torno de 47 \% e com o intervalo entre os extremos inferior a 5 \% ; variando de 44,1 \% no 1º quarto a 48,9 \% no último quarto.
No Gráfico 4.6, pode -se constatar que para os quatro quartos de desempenho a proporção de estudantes que consideraram a prova longa ou muito longa, em relação ao tempo total destinado à sua resolução, também apresentou percentuais próximos e manteve -se em torno de 47 \%, sem diferenças estatisticamente significativas. 4.3 COMPREENSÃO DOS ENUNCIADOS DAS QUESTÕES 4.3.1 Componente de Formação Geral Com relação aos enunciados das questões do Componente de Formação Geral ( Questão 4 ), as opiniões foram positivas, já que 74,2 \% dos alunos avaliados consideraram os enunciados de todas ou da maioria das questões claros e objetivos ( Gráfico 4.7, Gráfico 4.8, e, no Anexo II, a Tabela II.4 ).
Na análise regional, a percentagem de estudantes que avaliaram que todos ou a maioria dos enunciados das questões do Componente de Formação Geral estavam claros e objetivos variou de 71,7 \% na região Nordeste a 76,8 \% na região Sul, não havendo diferenças estatisticamente significativas entre as Grandes Regiões.
A análise das percepções dos estudantes sobre a clareza e objetividade dos enunciados permite afirmar que todos ou a maioria dos enunciados de questões relativas ao Componente de Formação Geral foram considerados claros e objetivos para a maior parte dos respondentes ( maior do que 71 \% em todas as regiões e maior do que 61 \% para todos os quartos de desempenho ).
Segundo o desempenho, observa -se que a proporção dos que emitiram esta opinião cresce conforme o desempenho aumenta, com diferenças estatisticamente significativas.
No quarto superior, a clareza e objetividade de todos ou da maioria dos enunciados das questões foi percebida por 79,6 \% dos alunos e no quarto de desempenho inferior tal avaliação foi emitida por 61,3 \% deles, numa proporção crescente com o desempenho nos três primeiros quartos. 4.3.2 Componente de Conhecimento Específico Com relação aos enunciados das questões do Componente de Conhecimento Específico da prova, para 67,4 \% dos estudantes avaliados da Área de Filosofia a clareza e a objetividade ( Questão 5 ) estavam presentes em todas ou na maioria das questões ( Gráfico 4.9, Gráfico 4.10, e no Anexo II, a Tabela II.5 ).
A maioria dos estudantes de todas as Grandes Regiões brasileiras considerou claros e objetivos todos ou a maioria dos enunciados das questões do Componente de Conhecimento Específico da prova, percentual sempre maior do que 60 \%.
A diferença entre 63,8 \% ( Nordeste ) e 69,6 \% ( Norte ) não é estatisticamente significativa.
A proporção de estudantes que consideraram os enunciados das questões claros e objetivos apresenta uma tendência crescente em relação ao aumento de desempenho : mais elevada no 3º quarto ( 72,9 \% ) se comparada ao quarto inferior de desempenho ( 56,3 \% ).
Esta diferença é estatisticamente significativa. 4.4 SUFICIÊNCIA DAS INFORMAÇÕES/INSTRUÇÕES FORNECIDAS Ao avaliarem as informações/instruções fornecidas para a resolução das questões ( Questão 6 ), 78,4 \% dos respondentes da Área de Filosofia de todo o Brasil afirmaram que estas eram até excessivas ou suficientes em todas ou na maioria das questões ( Gráfico 4.11, Gráfico 4.12, e, no Anexo II, a Tabela II.6 ).
Quanto à distribuição de respondentes pelas Grandes Regiões observa -se que a proporção de estudantes que consideraram as informações/instruções fornecidas até excessivas ou suficientes em todas ou na maioria das questões foi sempre superior a 70 \%, chegando a 82,6 \% na região Sul.
As diferenças entre as regiões não chegaram a ser estatisticamente significativa, a exceção da comparação entre a região Nordeste e a região Sul.
Levando -se em conta o desempenho dos participantes, notam -se diferenças estatisticamente significativas entre as opiniões de estudantes do quarto inferior e superior de desempenho, como mostra o Gráfico 4.12.
O percentual de participantes que avaliaram as informações/instruções como até excessivas ou suficientes em todas ou na maioria das questões foi mais elevado no quarto superior ( 81,7 \% ), percentual superior à média nacional ( 78,4 \% ).
Já no quarto inferior, a suficiência das informações/instruções declarada como até excessiva em todas ou na maioria das questões foi percebida por 70,5 \% dos respondentes.
Nota -se uma tendência crescente com a desempenho na proporção dos alunos que escolhem esta alternativa. 4.5 DIFICULDADE ENCONTRADA AO RESPONDER À PROVA Perguntados sobre as dificuldades com as quais se depararam ao responder à prova ( Questão 7 ), 16,0 \% dos estudantes apontaram o desconhecimento do conteúdo.
Para 45,5 \%, a forma diferente de abordagem do conteúdo foi indicada como dificuldade.
Já a falta de motivação para fazer a prova foi a dificuldade apontada por 21,8 \% dos respondentes.
Considerando -se todo o Brasil, 11,2 \% dos respondentes afirmaram que não tiveram qualquer tipo de dificuldade para responder à prova ( Tabela II.7 no Anexo II ).
Os Gráficos 4.13 e 4.14 apresentam os percentuais de estudantes que apontaram o desconhecimento do conteúdo como dificuldade percebida ao responder à prova.
Na análise por Grandes Regiões, o percentual de inscritos e presentes que apontaram o desconhecimento do conteúdo como dificuldade ao responder à prova não superou 18 \%.
Os percentuais variaram de 10,7 \% na região Centro-Oeste a 17,3 \% nas regiões Norte e Sudeste, sendo a diferença entre as regiões Centro-Oeste e Sudeste significativa estatisticamente.
A forma diferente de abordagem do conteúdo foi a escolha modal dos estudantes, com percentuais que variaram de 38,6 \% ( região Sudeste ) a 54,9 \% ( Centro-Oeste ).
O percentual de alunos que citou a falta de motivação como dificuldade variou de 17,2 \% ( região Centro-Oeste ) a 25,1 \% ( região Sudeste ).
Os que declararam não ter qualquer dificuldade para responder à prova variaram de 7,8 \% na região Nordeste a 14,2 \% na Centro-Oeste.
Com relação aos quartos de desempenho, o desconhecimento do conteúdo, foi a opção escolhida por 15,9 \% dos estudantes do quarto superior e 15,5 \% do quarto inferior.
A alternativa modal para os alunos, quando agregados pelos quartos de desempenho, foi que a dificuldade encontrada foi causada pela forma diferente de abordagem do conteúdo : 40,1 \% no quarto inferior e 45,9 \% do quarto superior assim o responderam. 4.6 CONTEÚDOS DAS QUESTÕES OBJETIVAS DA PROVA Ao analisarem os conteúdos das questões objetivas da prova ( Questão 8 ), um percentual pequeno dos estudantes avaliados, apenas 8,7 \%, afirmou que não estudou ainda a maioria desses conteúdos ( Gráficos 4.15, Gráfico 4.16, e a Tabelas II.8 no Anexo II ).
A maioria ( 59,2 \% ) afirmou ter estudado e aprendido muitos ou todos os conteúdos.
Na análise por Grande Região, a proporção de respondentes que escolheram a opção não estudou ainda a maioria desses conteúdos, foi pequena.
Observa -se que nas regiões Norte ( 10,0 \% ) e Nordeste ( 12,1 \% ), apesar de pequenas, as proporções foram maiores do que a média nacional ( 8,7 \% ).
Observa -se diferença estatisticamente significativa entre as regiões Sul e Nordeste.
Em todas as regiões, a maioria dos presentes afirmou ter estudado e aprendido muitos ou todos os conteúdos, com proporções variando entre 49,4 \% na região Nordeste e 64,7 \% na Centro-Oeste.
Considerando -se separadamente as opiniões de estudantes dos quatro quartos de desempenho, observa -se que, no quarto inferior, 14,4 \% ofereceram como resposta que não estudou ainda a maioria desses conteúdos, sendo menos de 4,5 \% os do quarto superior com a mesma resposta.
A diferença entre os alunos que optaram por este motivo de dificuldade nos quartos extremos é estatisticamente significativa e as proporções evidenciam uma tendência decrescente com o desempenho.
Tendo em conta o quarto superior, 75,1 \% dos alunos afirmaram ter estudado e aprendido muitos ou todos os conteúdos. 4.7 TEMPO GASTO PARA CONCLUIR A PROVA Ao responderem sobre o tempo de conclusão da prova ( Questão 9 ), mais de dois terços dos estudantes ( 71,0 \% ) afirmaram ter gasto entre duas e quatro horas ( Gráfico 4.17, Gráfico 4.18 e, no Anexo II, a Tabela II.9 ).
Considerando -se as cinco Grandes Regiões brasileiras, o percentual dos que utilizaram entre duas e quatro horas para finalizar a prova ficou abaixo do percentual nacional na região Sudeste ( 68,0 \% ).
Nas demais regiões o percentual de alunos que dispensaram entre duas e quatro horas para concluir a prova, ficou em torno de 74 \%, como mostra o Gráfico 4.17.
Uma vez sendo analisadas as alternativas escolhidas pelos estudantes que se situam nos diferentes quartos de desempenho, observa -se que uma maior proporção de participantes no quarto superior declarou ter gasto entre duas e quatro horas para concluir a prova quando comparados com os do quarto inferior, respectivamente 79,6 \% e 55,1 \%.
As diferenças entre os quartos extremos são estatisticamente significativas e os valores evidenciam uma tendência crescente com o desempenho.

% matched lemmas: 
\end{textsample}
